%==============================================================================
% PRESENTACIÓN: DINÁMICA MOLECULAR
% Universidad de Caldas - Departamento de Química
% Introducción a la Química Computacional (173G7G)
% Profesor: José Mauricio Rodas Rodríguez
%==============================================================================

\documentclass[aspectratio=169, 10pt]{beamer}

%------------------------------------------------------------------------------
% PAQUETES
%------------------------------------------------------------------------------
\usepackage[utf8]{inputenc}
\usepackage[T1]{fontenc}
\usepackage[spanish]{babel}
\usepackage{amsmath, amssymb, amsfonts}
\usepackage{mathtools}
\usepackage{physics}
\usepackage{chemformula}
\usepackage{booktabs}
\usepackage{multirow}
\usepackage{array}
\usepackage{colortbl}
\usepackage{xcolor}
\usepackage{tikz}
\usepackage{pgfplots}
\usepackage{pgfplotstable}
\usepackage{fontawesome5}
\usepackage{tcolorbox}
\usepackage{enumerate}
\usepackage{graphicx}
\usepackage{lmodern}
\usepackage{microtype}
\usepackage{hyperref}
\usepackage{listings}

\pgfplotsset{compat=1.18}
\usetikzlibrary{
  arrows.meta,
  shapes.geometric,
  shapes.misc,
  positioning,
  fit,
  backgrounds,
  decorations.pathmorphing,
  decorations.markings,
  calc,
  shadows,
  patterns
}
\tcbuselibrary{skins, breakable, theorems}

%------------------------------------------------------------------------------
% PALETA DE COLORES PASTEL — MODERNA
%------------------------------------------------------------------------------
\definecolor{PastelBlue}{RGB}{174, 214, 241}       % Azul suave
\definecolor{PastelGreen}{RGB}{171, 235, 198}      % Verde suave
\definecolor{PastelPurple}{RGB}{212, 172, 220}     % Lila
\definecolor{PastelOrange}{RGB}{250, 215, 160}     % Naranja suave
\definecolor{PastelRed}{RGB}{245, 183, 177}        % Rojo-rosa
\definecolor{PastelYellow}{RGB}{249, 231, 159}     % Amarillo suave
\definecolor{PastelTeal}{RGB}{163, 228, 215}       % Teal
\definecolor{DarkBlue}{RGB}{31, 97, 141}           % Azul oscuro (texto)
\definecolor{DarkGreen}{RGB}{30, 132, 73}          % Verde oscuro
\definecolor{DarkPurple}{RGB}{118, 68, 138}        % Morado oscuro
\definecolor{AccentGray}{RGB}{245, 246, 250}       % Fondo claro
\definecolor{TitleColor}{RGB}{21, 67, 96}          % Título oscuro
\definecolor{SubtitleColor}{RGB}{93, 109, 126}     % Subtítulo

%------------------------------------------------------------------------------
% TEMA BEAMER PERSONALIZADO
%------------------------------------------------------------------------------
\usetheme{Madrid}
\usecolortheme{default}

% Estructura de colores
\setbeamercolor{palette primary}{bg=DarkBlue, fg=white}
\setbeamercolor{palette secondary}{bg=PastelBlue, fg=DarkBlue}
\setbeamercolor{palette tertiary}{bg=DarkBlue!80, fg=white}
\setbeamercolor{palette quaternary}{bg=DarkBlue, fg=white}

\setbeamercolor{structure}{fg=DarkBlue}
\setbeamercolor{frametitle}{bg=DarkBlue, fg=white}
\setbeamercolor{frametitle right}{bg=PastelBlue}
\setbeamercolor{title}{fg=white}
\setbeamercolor{subtitle}{fg=white}
\setbeamercolor{author}{fg=DarkBlue}
\setbeamercolor{institute}{fg=SubtitleColor}
\setbeamercolor{date}{fg=SubtitleColor}
\setbeamercolor{background canvas}{bg=white}

\setbeamercolor{block title}{bg=DarkBlue, fg=white}
\setbeamercolor{block body}{bg=PastelBlue!30, fg=black}
\setbeamercolor{block title alerted}{bg=DarkPurple, fg=white}
\setbeamercolor{block body alerted}{bg=PastelPurple!40, fg=black}
\setbeamercolor{block title example}{bg=DarkGreen, fg=white}
\setbeamercolor{block body example}{bg=PastelGreen!40, fg=black}

\setbeamercolor{itemize item}{fg=DarkBlue}
\setbeamercolor{itemize subitem}{fg=DarkPurple}
\setbeamercolor{enumerate item}{fg=DarkBlue}

% Fuentes
\setbeamerfont{title}{size=\Large, series=\bfseries}
\setbeamerfont{subtitle}{size=\normalsize}
\setbeamerfont{frametitle}{size=\large, series=\bfseries}
\setbeamerfont{block title}{size=\normalsize, series=\bfseries}

% Eliminar barra de navegación
\setbeamertemplate{navigation symbols}{}

% Pie de página personalizado
\setbeamertemplate{footline}{
  \leavevmode
  \hbox{%
    \begin{beamercolorbox}[wd=0.33\paperwidth, ht=2.5ex, dp=1.5ex, leftskip=0.3cm]{palette primary}%
      \usebeamerfont{author in head/foot}\insertshortauthor
    \end{beamercolorbox}%
    \begin{beamercolorbox}[wd=0.34\paperwidth, ht=2.5ex, dp=1.5ex, center]{palette secondary}%
      \usebeamerfont{title in head/foot}\insertshorttitle
    \end{beamercolorbox}%
    \begin{beamercolorbox}[wd=0.33\paperwidth, ht=2.5ex, dp=1.5ex, rightskip=0.3cm, right]{palette primary}%
      \usebeamerfont{date in head/foot}
      \insertshortdate{}\hspace*{2em}
      \insertframenumber{}/\inserttotalframenumber
    \end{beamercolorbox}%
  }%
}

% Encabezado de sección
\setbeamertemplate{headline}{
  \begin{beamercolorbox}[ht=3pt]{palette primary}
  \end{beamercolorbox}
}

%------------------------------------------------------------------------------
% CAJAS PERSONALIZADAS (tcolorbox)
%------------------------------------------------------------------------------
\tcbset{
  formulabox/.style={
    enhanced, breakable,
    colback=PastelBlue!25,
    colframe=DarkBlue,
    fonttitle=\bfseries,
    boxrule=1.5pt,
    arc=4pt,
    left=6pt, right=6pt, top=4pt, bottom=4pt
  },
  definitionbox/.style={
    enhanced, breakable,
    colback=PastelGreen!30,
    colframe=DarkGreen,
    fonttitle=\bfseries,
    boxrule=1.5pt,
    arc=4pt,
    left=6pt, right=6pt, top=4pt, bottom=4pt
  },
  warningbox/.style={
    enhanced, breakable,
    colback=PastelOrange!40,
    colframe=DarkPurple!80,
    fonttitle=\bfseries,
    boxrule=1.5pt,
    arc=4pt,
    left=6pt, right=6pt, top=4pt, bottom=4pt
  },
  exercisebox/.style={
    enhanced, breakable,
    colback=PastelYellow!40,
    colframe=DarkPurple,
    fonttitle=\bfseries,
    boxrule=1.5pt,
    arc=4pt,
    left=6pt, right=6pt, top=4pt, bottom=4pt
  },
  comparisonbox/.style={
    enhanced,
    colback=AccentGray,
    colframe=DarkBlue!60,
    boxrule=1pt,
    arc=4pt,
    left=4pt, right=4pt, top=3pt, bottom=3pt
  }
}

%------------------------------------------------------------------------------
% ESTILOS DE CÓDIGO (listings)
%------------------------------------------------------------------------------
\lstdefinelanguage{mdpfile}{%
  morecomment=[l]{;}%
}
\lstdefinestyle{bashstyle}{%
  language=bash,
  basicstyle=\ttfamily\tiny,
  keywordstyle=\color{DarkBlue}\bfseries,
  commentstyle=\color{DarkGreen!70}\itshape,
  stringstyle=\color{DarkPurple},
  backgroundcolor=\color{AccentGray},
  frame=single, framerule=0.4pt, rulecolor=\color{PastelBlue},
  breaklines=true, breakatwhitespace=true,
  showstringspaces=false, tabsize=2,
  xleftmargin=3pt, xrightmargin=3pt
}
\lstdefinestyle{mdpstyle}{%
  language=mdpfile,
  basicstyle=\ttfamily\tiny,
  commentstyle=\color{DarkGreen!70}\itshape,
  backgroundcolor=\color{PastelBlue!15},
  frame=single, framerule=0.4pt, rulecolor=\color{PastelBlue},
  breaklines=true, showstringspaces=false,
  tabsize=2, xleftmargin=3pt, xrightmargin=3pt
}

%------------------------------------------------------------------------------
% INFORMACIÓN DEL DOCUMENTO
%------------------------------------------------------------------------------
\title[Dinámica Molecular]{Dinámica Molecular}
\subtitle{Principios, Métodos y Aplicaciones en Química Computacional}
\author[J.M. Rodas R.]{Prof. José Mauricio Rodas Rodríguez}
\institute[U. Caldas]{Universidad de Caldas \\ Departamento de Química \\ \textit{Introducción a la Química Computacional (173G7G)}}
\date{\today}

%==============================================================================
\begin{document}
%==============================================================================

%------------------------------------------------------------------------------
% DIAPOSITIVA TÍTULO
%------------------------------------------------------------------------------
\begin{frame}[plain]
\titlepage
\vspace{-1em}
\begin{center}
  \begin{tikzpicture}[scale=0.8]
    % Representación esquemática de partículas en MD
    \foreach \x/\y/\c in {
      0/0/PastelBlue, 1.2/0.3/PastelGreen, 0.5/1/PastelPurple,
      2/0.8/PastelOrange, 1.5/1.5/PastelTeal, 2.8/0.2/PastelRed,
      3.2/1.2/PastelYellow, 0.2/2/PastelBlue!60}{
      \node[circle, fill=\c, draw=DarkBlue!40, minimum size=14pt] at (\x,\y) {};
    }
    % Flechas de velocidad
    \draw[->, thick, DarkBlue!70] (0,0) -- (0.4, 0.3);
    \draw[->, thick, DarkGreen!70] (1.2,0.3) -- (0.8, 0.7);
    \draw[->, thick, DarkPurple!70] (0.5,1) -- (0.9, 1.3);
    \draw[->, thick, DarkBlue!50] (2,0.8) -- (2.5, 0.6);
  \end{tikzpicture}
\end{center}
\end{frame}

%------------------------------------------------------------------------------
% TABLA DE CONTENIDOS
%------------------------------------------------------------------------------
\begin{frame}{Contenido}
\begin{columns}[T]
\column{0.5\textwidth}
\tableofcontents[hideallsubsections, sections={1-5}]
\column{0.5\textwidth}
\tableofcontents[hideallsubsections, sections={6-9}]
\end{columns}
\end{frame}

%==============================================================================
\section{Fundamentos de Dinámica Molecular}
%==============================================================================

\begin{frame}{¿Qué es la Dinámica Molecular?}
\begin{columns}[T]
\column{0.55\textwidth}
\begin{tcolorbox}[definitionbox, title={Definición}]
La \textbf{Dinámica Molecular (MD)} es una técnica de simulación computacional que resuelve las ecuaciones de movimiento de Newton para un conjunto de átomos o moléculas en función del tiempo.
\end{tcolorbox}
\vspace{0.5em}
\begin{itemize}
  \item Permite estudiar el \textbf{comportamiento dinámico} de sistemas a escala atómica
  \item Proporciona información estructural, termodinámica y cinética
  \item Escala de tiempo: fs a $\mu$s; escala espacial: Å a nm
\end{itemize}

\column{0.42\textwidth}
\begin{tcolorbox}[formulabox, title={Ecuación de Newton}]
\[
m_i \ddot{\mathbf{r}}_i = \mathbf{F}_i = -\nabla_i U(\mathbf{r}^N)
\]
\small
donde $U(\mathbf{r}^N)$ es la energía potencial del sistema de $N$ partículas.
\end{tcolorbox}
\vspace{0.3em}
\begin{tcolorbox}[warningbox, title={Hipótesis clave}]
\small
\begin{itemize}
  \item Núcleos tratados clásicamente
  \item Electrones implícitos (campo de fuerzas)
  \item Superficie de Born-Oppenheimer
\end{itemize}
\end{tcolorbox}
\end{columns}
\end{frame}

%------------------------------------------------------------------------------
\begin{frame}[shrink=15]{Campo de Fuerzas: Energía Potencial}
\begin{tcolorbox}[formulabox, title={Función de Energía Potencial Total}]
\scriptsize
\begin{align*}
U &= \underbrace{\sum_{\text{enlaces}} \frac{k_b}{2}(b - b_0)^2}_{\text{enlace}}
  + \underbrace{\sum_{\text{ángulos}} \frac{k_\theta}{2}(\theta - \theta_0)^2}_{\text{ángulo}}
  + \underbrace{\sum_{\text{diedros}} \frac{V_n}{2}[1+\cos(n\phi - \delta)]}_{\text{torsión}} \\
  &\quad + \underbrace{\sum_{i<j} 4\varepsilon_{ij}\!\left[\left(\frac{\sigma_{ij}}{r_{ij}}\right)^{12}\!-\!\left(\frac{\sigma_{ij}}{r_{ij}}\right)^{6}\right]}_{\text{Lennard-Jones}}
  + \underbrace{\sum_{i<j} \frac{q_i q_j}{4\pi\varepsilon_0 r_{ij}}}_{\text{Coulomb}}
\end{align*}
\end{tcolorbox}

\vspace{0.1em}
\begin{columns}[T]
\column{0.48\textwidth}
\begin{block}{Términos ligados (bonded)}
  \begin{itemize}
    \item Estiramiento de enlace ($k_b$, $b_0$)
    \item Flexión de ángulo ($k_\theta$, $\theta_0$)
    \item Torsión y diedros impropios ($V_n$, $\phi$, $\delta$)
  \end{itemize}
\end{block}

\column{0.48\textwidth}
\begin{block}{Términos no ligados (non-bonded)}
  \begin{itemize}
    \item Van der Waals: potencial LJ ($\varepsilon$, $\sigma$)
    \item Electrostático: cargas parciales ($q_i$)
    \item Cutoff o suma de Ewald
  \end{itemize}
\end{block}
\end{columns}
\end{frame}

%------------------------------------------------------------------------------
\begin{frame}{Potencial de Lennard-Jones}
\begin{columns}[T]
\column{0.44\textwidth}
\begin{tcolorbox}[formulabox, title={Potencial LJ 12-6}]
\[
U_{LJ}(r) = 4\varepsilon\left[\left(\frac{\sigma}{r}\right)^{12} - \left(\frac{\sigma}{r}\right)^{6}\right]
\]
\small
\begin{itemize}
  \item $\varepsilon$: profundidad del pozo de energía
  \item $\sigma$: distancia donde $U_{LJ}=0$
  \item $r_{min} = 2^{1/6}\sigma$: mínimo del potencial
\end{itemize}
\end{tcolorbox}

\column{0.53\textwidth}
\begin{tikzpicture}
\begin{axis}[
  width=7cm, height=5.5cm,
  xlabel={$r/\sigma$},
  ylabel={$U/\varepsilon$},
  xmin=0.85, xmax=2.8,
  ymin=-1.5, ymax=3,
  ymajorgrids=true, xmajorgrids=true,
  grid style={gray!20},
  axis lines=left,
  tick label style={font=\tiny},
  label style={font=\small},
  title style={font=\small},
  title={Potencial de Lennard-Jones}
]
\addplot[thick, color=DarkBlue, domain=0.9:2.8, samples=100]{4*(x^(-12) - x^(-6))};
\addplot[dashed, color=DarkPurple, domain=0.9:2.8]{0};
\addplot[dotted, color=DarkGreen, domain=0.9:2.8]{-1};
\node[font=\tiny, color=DarkGreen] at (axis cs:2.3,-1.2) {$-\varepsilon$};
\node[font=\tiny, color=DarkBlue] at (axis cs:1.3,2.5) {$U_{LJ}$};
\end{axis}
\end{tikzpicture}
\end{columns}
\end{frame}

%------------------------------------------------------------------------------
\begin{frame}{Distribución de Maxwell-Boltzmann}
\begin{columns}[T]
\column{0.5\textwidth}
\begin{tcolorbox}[formulabox, title={Distribución de velocidades}]
\[
f(v) = 4\pi \left(\frac{m}{2\pi k_B T}\right)^{3/2} v^2 \exp\!\left(-\frac{mv^2}{2k_BT}\right)
\]
\end{tcolorbox}
\vspace{0.3em}
\begin{itemize}
  \item Velocidad más probable: $v_p = \sqrt{\dfrac{2k_BT}{m}}$
  \item Velocidad media: $\bar{v} = \sqrt{\dfrac{8k_BT}{\pi m}}$
  \item Velocidad cuadrática media: $v_{rms} = \sqrt{\dfrac{3k_BT}{m}}$
  \item Temperatura: $T = \dfrac{2\langle E_k \rangle}{3 N k_B}$
\end{itemize}

\column{0.50\textwidth}
\begin{tikzpicture}
\begin{axis}[
  width=7cm, height=5.5cm,
  xlabel={Velocidad (m/s)},
  ylabel={$f(v)$ (u.a.)},
  xmin=0, xmax=1000,
  ymin=0,
  ymajorgrids=true, xmajorgrids=true,
  grid style={gray!20},
  axis lines=left,
  tick label style={font=\tiny},
  label style={font=\small},
  title={Maxwell-Boltzmann ($N_2$)},
  title style={font=\small},
  legend style={font=\tiny, at={(0.98,0.98)}, anchor=north east}
]
\addplot[thick, DarkBlue, domain=0:1000, samples=200]{4*3.14159*(28.014e-3/(2*3.14159*8.314*300))^(1.5) * x^2 * exp(-28.014e-3*x^2/(2*8.314*300)) * 1e8};
\addlegendentry{300 K}
\addplot[thick, DarkPurple, dashed, domain=0:1000, samples=200]{4*3.14159*(28.014e-3/(2*3.14159*8.314*700))^(1.5) * x^2 * exp(-28.014e-3*x^2/(2*8.314*700)) * 1e8};
\addlegendentry{700 K}
\end{axis}
\end{tikzpicture}
\end{columns}
\end{frame}

%==============================================================================
\section{Algoritmos de Integración}
%==============================================================================

\begin{frame}{Algoritmo de Verlet}
\begin{columns}[T]
\column{0.52\textwidth}
\begin{tcolorbox}[formulabox, title={Expansión de Taylor}]
\small
\begin{align*}
\mathbf{r}(t+\Delta t) &= \mathbf{r}(t) + \mathbf{v}(t)\Delta t + \frac{\mathbf{F}(t)}{2m}\Delta t^2 \\
\mathbf{r}(t-\Delta t) &= \mathbf{r}(t) - \mathbf{v}(t)\Delta t + \frac{\mathbf{F}(t)}{2m}\Delta t^2
\end{align*}
\textbf{Sumando:}
\[
\mathbf{r}(t+\Delta t) = 2\mathbf{r}(t) - \mathbf{r}(t-\Delta t) + \frac{\mathbf{F}(t)}{m}\Delta t^2
\]
\textbf{Velocidades:}
\[
\mathbf{v}(t) = \frac{\mathbf{r}(t+\Delta t) - \mathbf{r}(t-\Delta t)}{2\Delta t}
\]
\end{tcolorbox}

\column{0.45\textwidth}
\begin{tcolorbox}[warningbox, title={Propiedades}]
\small
\begin{itemize}
  \item Error local: $\mathcal{O}(\Delta t^4)$
  \item Error global: $\mathcal{O}(\Delta t^2)$
  \item \textbf{Simpléctico} (conserva volumen de fase)
  \item Reversible en el tiempo
  \item No almacena velocidades explícitamente
\end{itemize}
\end{tcolorbox}
\vspace{0.3em}
\begin{tcolorbox}[definitionbox, title={Paso de tiempo típico}]
\small
$\Delta t \sim 1\text{--}2$ fs para \ch{H} \newline
$\Delta t \sim 2$ fs con restricciones \newline
$\Delta t \sim 4\text{--}8$ fs con H-virtual
\end{tcolorbox}
\end{columns}
\end{frame}

%------------------------------------------------------------------------------
\begin{frame}[shrink=25]{Velocity-Verlet y Leap-Frog}
\begin{columns}[T]
\column{0.48\textwidth}
\begin{tcolorbox}[formulabox, title={Velocity-Verlet \normalfont\scriptsize(3 pasos)}, equal height group=ehg-vv-eq]
\small
\begin{align*}
\mathbf{v}(t+\tfrac{\Delta t}{2}) &= \textcolor{DarkPurple}{\mathbf{v}(t)}
    + \frac{\mathbf{F}(t)\,\Delta t}{2m} \\
\mathbf{r}(t+\Delta t) &= \mathbf{r}(t) + \mathbf{v}(t+\tfrac{\Delta t}{2})\,\Delta t \\
\mathbf{v}(t+\Delta t) &= \mathbf{v}(t+\tfrac{\Delta t}{2})
    + \frac{\mathbf{F}(t+\Delta t)\,\Delta t}{2m}
\end{align*}
\end{tcolorbox}
\vspace{0.2em}
\begin{tcolorbox}[definitionbox, title={Propiedades}, equal height group=ehg-vv-pr]
\small
\begin{itemize}
  \item Parte de \textcolor{DarkPurple}{$\mathbf{v}(t)$} en paso \textbf{entero}
  \item $\mathbf{r}(t)$ y $\mathbf{v}(t)$ disponibles \textbf{simultáneamente}
  \item Evalúa $\mathbf{F}$ \textbf{dos veces} por paso ($t$ y $t{+}\Delta t$)
  \item GROMACS \texttt{md-vv}, NAMD
\end{itemize}
\end{tcolorbox}

\column{0.48\textwidth}
\begin{tcolorbox}[warningbox, title={Leap-Frog \normalfont\scriptsize(2 pasos)}, equal height group=ehg-vv-eq]
\small
\begin{align*}
\mathbf{v}(t+\tfrac{\Delta t}{2}) &= \textcolor{DarkGreen}{\mathbf{v}(t-\tfrac{\Delta t}{2})}
    + \frac{\mathbf{F}(t)\,\Delta t}{m} \\[4pt]
\mathbf{r}(t+\Delta t) &= \mathbf{r}(t) + \mathbf{v}(t+\tfrac{\Delta t}{2})\,\Delta t
\end{align*}
\end{tcolorbox}
\vspace{0.2em}
\begin{tcolorbox}[definitionbox, title={Propiedades}, equal height group=ehg-vv-pr]
\small
\begin{itemize}
  \item Parte de \textcolor{DarkGreen}{$\mathbf{v}(t-\tfrac{\Delta t}{2})$} en paso \textbf{medio}
  \item $\mathbf{r}(t)$ y $\mathbf{v}(t)$ en \textbf{instantes distintos} ($\Delta t/2$)
  \item Evalúa $\mathbf{F}$ \textbf{una sola vez} por paso (más eficiente)
  \item Estándar GROMACS \texttt{md}
\end{itemize}
\end{tcolorbox}
\end{columns}

\vspace{0.1em}
% Diagrama de línea temporal
\begin{center}
\begin{tikzpicture}[scale=0.72]
  % eje de tiempo
  \draw[->,gray!50,thin] (1.7,0)--(11.4,0)
    node[right,font=\tiny,gray]{tiempo};
  % marcas en pasos enteros
  \foreach \x/\lab in {2.5/{t-\Delta t},6.5/{t},10.5/{t+\Delta t}}{
    \draw[gray!50,thin](\x,-.05)--(\x,.05);
    \node[font=\tiny,gray,below=1pt] at (\x,0){$\lab$};
  }
  % marcas en pasos medios (punteadas)
  \foreach \x/\lab in {4.5/{t-\frac{\Delta t}{2}},8.5/{t+\frac{\Delta t}{2}}}{
    \draw[gray!40,dashed,thin](\x,-.05)--(\x,.05);
    \node[font=\tiny,gray!60,below=1pt] at (\x,0){$\lab$};
  }
  % VV: r y v en pasos enteros alineados
  \node[left,font=\tiny\bfseries,DarkBlue] at (1.7,.55){VV};
  \draw[DarkBlue!40,thin](2.5,.40)--(10.5,.40);
  \draw[DarkPurple!50,thin](2.5,.65)--(10.5,.65);
  \foreach \x in {2.5,6.5,10.5}{
    \fill[DarkBlue](\x,.40) circle(2pt);
    \fill[DarkPurple](\x,.65) circle(2pt);
  }
  \node[right,font=\tiny,DarkBlue] at (10.65,.40){$\mathbf{r}$};
  \node[right,font=\tiny,DarkPurple] at (10.65,.65){$\mathbf{v}$};
  % LF: r en pasos enteros, v en pasos medios (desfasado)
  \node[left,font=\tiny\bfseries,DarkGreen] at (1.7,-.55){LF};
  \draw[DarkBlue!40,thin](2.5,-.40)--(10.5,-.40);
  \draw[DarkGreen!60,thin](4.5,-.65)--(8.5,-.65);
  \foreach \x in {2.5,6.5,10.5}{\fill[DarkBlue](\x,-.40) circle(2pt);}
  \foreach \x in {4.5,8.5}{\fill[DarkGreen](\x,-.65) circle(2pt);}
  \node[right,font=\tiny,DarkBlue] at (10.65,-.40){$\mathbf{r}$};
  \draw[->,DarkGreen!70,thin](4.5,-.65) to[bend left=20]
    node[above,font=\tiny,DarkGreen]{salto $\Delta t$}(8.5,-.65);
\end{tikzpicture}
\end{center}
\end{frame}

%------------------------------------------------------------------------------
\begin{frame}{Flujo de Trabajo de un Paso MD}

\begin{center}
\begin{tikzpicture}[
  node distance=0.5cm,
  process/.style={rectangle, rounded corners=5pt, minimum width=3.2cm, minimum height=0.6cm,
    text centered, font=\small, draw=DarkBlue, fill=PastelBlue!50},
  decision/.style={diamond, minimum width=2.2cm, minimum height=0.5cm,
    text centered, font=\small, draw=DarkPurple, fill=PastelPurple!40, aspect=3},
  arrow/.style={->, thick, DarkBlue}
]

\node[process, fill=PastelGreen!60] (init) {Posiciones $\mathbf{r}(t)$, velocidades $\mathbf{v}(t)$};
\node[process, below=of init] (forces) {Calcular fuerzas $\mathbf{F}_i = -\nabla_i U$};
\node[process, below=of forces] (integrate) {Integrar: $\mathbf{r}(t{+}\Delta t)$, $\mathbf{v}(t{+}\Delta t)$};
\node[process, below=of integrate] (thermo) {Aplicar termostato/barostato};
\node[decision, below=of thermo] (check) {¿$t \geq t_{max}$?};
\node[process, right=2.8cm of check, fill=PastelOrange!50] (output) {Guardar trayectoria};
\node[process, below=of check, fill=PastelGreen!60] (end) {Análisis y resultados};

\draw[arrow] (init) -- (forces);
\draw[arrow] (forces) -- (integrate);
\draw[arrow] (integrate) -- (thermo);
\draw[arrow] (thermo) -- (check);
\draw[arrow] (check) -- node[right, font=\scriptsize]{Sí} (end);
\draw[arrow] (check) -- node[above, font=\scriptsize]{Guardar} (output);
\draw[arrow] (output) |- (forces);
\draw[arrow] (check.west) -- ++(-1.0,0) |- node[left, font=\scriptsize, near start]{No} (forces.west);

\end{tikzpicture}
\end{center}
\end{frame}

%==============================================================================
\section{Condiciones de Contorno y Ensambles}
%==============================================================================

\begin{frame}{Condiciones Periódicas de Contorno (PBC)}
\begin{columns}[T]
\column{0.52\textwidth}
\begin{block}{¿Por qué PBC?}
  \begin{itemize}
    \item Elimina efectos de borde artificiales
    \item Simula sistema "infinito" con $N$ finito
    \item Conserva la densidad correcta
    \item Permite calcular propiedades del bulto
  \end{itemize}
\end{block}
\vspace{0.3em}
\begin{tcolorbox}[formulabox, title={Convención de imagen mínima}]
\small
\[
\Delta x_{ij} = \Delta x_{ij} - L_x \cdot \text{round}\!\left(\frac{\Delta x_{ij}}{L_x}\right)
\]
La distancia más corta entre dos partículas se calcula usando la copia de imagen más cercana.
\end{tcolorbox}

\column{0.45\textwidth}
\begin{center}
\begin{tikzpicture}[scale=0.9]
% Celdas vecinas
\foreach \dx in {-1,0,1}{
  \foreach \dy in {-1,0,1}{
    \draw[gray!30, thin] (\dx*2.0, \dy*2.0) rectangle ++(2.0,2.0);
    \ifnum\dx=0
      \ifnum\dy=0
        \draw[DarkBlue, thick] (\dx*2.0, \dy*2.0) rectangle ++(2.0,2.0);
      \fi
    \fi
  }
}
% Partículas en la celda central
\foreach \px/\py in {0.4/0.5, 1.2/1.1, 0.8/1.7, 1.6/0.4}{
  \node[circle, fill=PastelBlue, draw=DarkBlue, minimum size=8pt] at (\px*2.0+0.0, \py*2.0-2.0+0.0) {};
}
% Imágenes
\node[circle, fill=PastelBlue!40, draw=DarkBlue!40, minimum size=8pt] at (0.4*2-2, 0.5*2-2+2) {};
\node[circle, fill=PastelBlue!40, draw=DarkBlue!40, minimum size=8pt] at (0.4*2+2, 0.5*2-2) {};
\node[font=\tiny, DarkBlue] at (1.0, -0.3) {Celda principal};
\end{tikzpicture}
\end{center}
\end{columns}
\end{frame}

%------------------------------------------------------------------------------
\begin{frame}{Ensambles Estadísticos}

\begin{columns}[T]
\column{0.32\textwidth}
\begin{tcolorbox}[formulabox, title={NVE (Microcanónico)}, equal height group=ehg-ens]
\small
\centering
$N$, $V$, $E$ = constantes \\[4pt]
\[
\mathcal{H} = K + U = E
\]
\vspace{2pt}
\begin{itemize}\small
  \item Sin intercambio de calor
  \item Dinámica newtoniana pura
  \item Validación de integradores
\end{itemize}
\end{tcolorbox}

\column{0.32\textwidth}
\begin{tcolorbox}[definitionbox, title={NVT (Canónico)}, equal height group=ehg-ens]
\small
\centering
$N$, $V$, $T$ = constantes \\[4pt]
\[
Z = \int e^{-\beta H}\,d\mathbf{r}\,d\mathbf{p}
\]
\vspace{2pt}
\begin{itemize}\small
  \item Termostato activo
  \item Función de partición canónica
  \item Más común en práctica
\end{itemize}
\end{tcolorbox}

\column{0.32\textwidth}
\begin{tcolorbox}[warningbox, title={NPT (Isotérmico-isob.)}, equal height group=ehg-ens]
\small
\centering
$N$, $P$, $T$ = constantes \\[4pt]
\[
\Delta G = \Delta H - T\Delta S
\]
\vspace{2pt}
\begin{itemize}\small
  \item Termostato + barostato
  \item Densidad variable
  \item Condiciones experimentales
\end{itemize}
\end{tcolorbox}
\end{columns}

\vspace{0.5em}
\begin{tcolorbox}[comparisonbox]
\centering\small
\textbf{Uso recomendado:} NVT para equilibración $\rightarrow$ NPT para producción en sistemas condensados
\end{tcolorbox}
\end{frame}

%==============================================================================
\section{Termostatos y Barostatos}
%==============================================================================

\begin{frame}{Termostatos}
\begin{columns}[T]
\column{0.48\textwidth}
\begin{tcolorbox}[formulabox, title={Berendsen}]
\small
\[
\frac{dT}{dt} = \frac{T_0 - T}{\tau_T}
\]
\[
\lambda = \sqrt{1 + \frac{\Delta t}{\tau_T}\left(\frac{T_0}{T} - 1\right)}
\]
{\scriptsize Escalado débil: no genera ensamble NVT correcto. Rápida estabilización.}
\end{tcolorbox}

\begin{tcolorbox}[definitionbox, title={Velocidad de rescalado (v-rescale)}]
\small
Añade término estocástico a Berendsen para recuperar ensamble canónico correcto. Recomendado en GROMACS.
\end{tcolorbox}

\column{0.48\textwidth}
\begin{tcolorbox}[warningbox, title={Nosé-Hoover}]
\small
Extiende el Hamiltoniano con variable ficticia $s$:
\[
\mathcal{H}_{NH} = \sum_i \frac{p_i^2}{2m_i s^2} + U + \frac{p_s^2}{2Q} + g k_B T \ln s
\]
Ecuaciones de movimiento extendidas:
\begin{align*}
\dot{\mathbf{r}}_i &= \frac{\mathbf{p}_i}{m_i} \\
\dot{\mathbf{p}}_i &= \mathbf{F}_i - \xi \mathbf{p}_i \\
\dot{\xi} &= \frac{1}{Q}\left(\sum_i \frac{p_i^2}{m_i} - g k_B T\right)
\end{align*}
\end{tcolorbox}
\end{columns}
\end{frame}

%------------------------------------------------------------------------------
\begin{frame}{Barostatos}
\begin{columns}[T]
\column{0.48\textwidth}
\begin{tcolorbox}[formulabox, title={Barostato de Berendsen}, equal height group=ehg-bar-a]
\small
\[
\frac{dP}{dt} = \frac{P_0 - P}{\tau_P}
\]
\[
\mu = \left[1 - \frac{\kappa_T \Delta t}{\tau_P}(P_0 - P)\right]^{1/3}
\]
Escala coordenadas y caja: $\mathbf{r} \rightarrow \mu\mathbf{r}$, $L \rightarrow \mu L$
\end{tcolorbox}
\vspace{0.3em}
\begin{itemize}
  \item $\kappa_T$: compresibilidad isotérmica
  \item Convergencia rápida al equilibrio
  \item No genera ensamble NPT correcto
\end{itemize}

\column{0.48\textwidth}
\begin{tcolorbox}[warningbox, title={Barostato de Parrinello-Rahman}, equal height group=ehg-bar-a]
\small
Generaliza la caja con tensor $\mathbf{h}$:
\[
\ddot{\mathbf{h}} = \frac{V}{W}\left(\mathbf{P} - P_0\mathbf{I}\right)\mathbf{h}^{-T}
\]
\begin{itemize}
  \item $W$: masa ficticia de la caja
  \item Permite cambio de forma y tamaño
  \item Genera ensamble NPT correcto
  \item Recomendado para producción
\end{itemize}
\end{tcolorbox}
\vspace{0.2em}
\begin{tcolorbox}[definitionbox, title={Secuencia típica}]
\scriptsize
Berendsen (equilibración) $\rightarrow$ Parrinello-Rahman (producción)
\end{tcolorbox}
\end{columns}
\end{frame}

%==============================================================================
\section{Preparación de Sistemas}
%==============================================================================

\begin{frame}{Protocolo de Preparación de una Simulación}

\begin{center}
\resizebox{0.98\textwidth}{!}{
\begin{tikzpicture}[
  node distance=0.35cm and 1.1cm,
  pstep/.style={rectangle, rounded corners=4pt, minimum width=2.5cm, minimum height=0.75cm,
    text width=2.4cm, align=center, font=\footnotesize, draw=DarkBlue, fill=PastelBlue!40},
  arrow/.style={->, thick, DarkBlue}
]
\node[pstep, fill=PastelGreen!60] (pdb) {1. Obtener estructura (PDB/AlphaFold)};
\node[pstep, right=of pdb] (clean) {2. Limpiar PDB (eliminar HETATM)};
\node[pstep, right=of clean] (ff) {3. Asignar campo de fuerzas};
\node[pstep, right=of ff] (box) {4. Definir caja de simulación};

\node[pstep, below=1.4cm of box] (solv) {5. Solvatación (añadir agua)};
\node[pstep, left=of solv] (ions) {6. Neutralizar (añadir iones)};
\node[pstep, left=of ions] (em) {7. Minimización de energía};
\node[pstep, left=of em, fill=PastelOrange!50] (eq) {8. Equilibración NVT / NPT};

\node[pstep, below=1.4cm of eq, fill=PastelYellow!60] (prod) {9. Producción (trayectoria)};
\node[pstep, right=of prod] (traj) {10. Guardar trayectoria};
\node[pstep, right=of traj, fill=PastelGreen!60] (analysis) {11. Análisis (RMSD, RMSF...)};

\draw[arrow] (pdb) -- (clean);
\draw[arrow] (clean) -- (ff);
\draw[arrow] (ff) -- (box);
\draw[arrow] (box) -- (solv);
\draw[arrow] (solv) -- (ions);
\draw[arrow] (ions) -- (em);
\draw[arrow] (em) -- (eq);
\draw[arrow] (eq) -- (prod);
\draw[arrow] (prod) -- (traj);
\draw[arrow] (traj) -- (analysis);
\end{tikzpicture}
}
\end{center}
\end{frame}

%------------------------------------------------------------------------------
\begin{frame}{Modelos de Agua}

\begin{columns}[T]
\column{0.55\textwidth}
\begin{table}[h]
\centering
\footnotesize
\begin{tabular}{@{}lcccc@{}}
\toprule
\textbf{Modelo} & \textbf{Sitios} & \textbf{Carga O} & \textbf{Ángulo} & \textbf{Uso} \\
\midrule
SPC      & 3 & $-0.82e$ & $109.47°$ & General \\
SPC/E    & 3 & $-0.8476e$ & $109.47°$ & \textbf{Estándar} \\
TIP3P    & 3 & $-0.834e$ & $104.52°$ & CHARMM \\
TIP4P    & 4 & $-1.04e$ & $104.52°$ & Mejorado \\
TIP4P/2005 & 4 & $-1.1128e$ & $104.52°$ & Mejor \\
TIP5P    & 5 & --- & $104.52°$ & Propiedades \\
OPC      & 4 & $-1.3582e$ & $103.6°$ & DNN \\
\bottomrule
\end{tabular}
\end{table}

\column{0.42\textwidth}
\begin{tcolorbox}[definitionbox, title={Recomendaciones}]
\small
\begin{itemize}
  \item \textbf{SPC/E}: uso general, buenas propiedades de difusión
  \item \textbf{TIP3P}: compatible con CHARMM/AMBER
  \item \textbf{TIP4P/2005}: mejores propiedades termofísicas
  \item \textbf{OPC}: alta precisión para biomoléculas
\end{itemize}
\end{tcolorbox}
\end{columns}
\end{frame}

%==============================================================================
\section{Campos de Fuerzas}
%==============================================================================

\begin{frame}{Principales Campos de Fuerzas}

\begin{columns}[T]
\column{0.48\textwidth}
\begin{block}{Para Proteínas y Biomoléculas}
  \begin{itemize}
    \item \textbf{AMBER}: ff14SB, ff19SB, GAFF (ligandos)
    \item \textbf{CHARMM}: CHARMM36, CGenFF
    \item \textbf{OPLS}: OPLS-AA, OPLS3e, OPLS4
    \item \textbf{GROMOS}: 54A7, 54A8
  \end{itemize}
\end{block}
\vspace{0.3em}
\begin{tcolorbox}[warningbox, title={Para moléculas pequeñas}, equal height group=ehg-ff-b]
\small
\begin{itemize}
  \item \textbf{GAFF2}: AMBER para ligandos
  \item \textbf{CGenFF}: CHARMM general
  \item \textbf{OPLS3e}: pequeñas moléculas
  \item \textbf{OpenFF}: Sage, Parsley
\end{itemize}
\end{tcolorbox}

\column{0.48\textwidth}
\begin{block}{Para Materiales y Polímeros}
  \begin{itemize}
    \item \textbf{DREIDING}: moléculas orgánicas
    \item \textbf{UFF}: fuerza universal (todos los átomos)
    \item \textbf{ReaxFF}: bonds reactivos
    \item \textbf{Lennard-Jones}: gases nobles
  \end{itemize}
\end{block}
\vspace{0.3em}
\begin{tcolorbox}[definitionbox, title={Validación crítica}, equal height group=ehg-ff-b]
\small
Un campo de fuerzas incorrecto produce \textbf{resultados erróneos}. Siempre verificar:
\begin{enumerate}
  \item Compatibilidad con el sistema
  \item Literatura de referencia
  \item Propiedades reproducidas
\end{enumerate}
\end{tcolorbox}
\end{columns}
\end{frame}

%==============================================================================
\section{Análisis de Trayectorias}
%==============================================================================

\begin{frame}{Métricas de Análisis: RMSD y RMSF}

\begin{columns}[T]
\column{0.48\textwidth}
\begin{tcolorbox}[formulabox, title={RMSD (Desviación Cuadrática Media)}, equal height group=ehg-rmsd]
\[
\text{RMSD}(t) = \sqrt{\frac{1}{N}\sum_{i=1}^{N} m_i \left|\mathbf{r}_i(t) - \mathbf{r}_i^{\text{ref}}\right|^2}
\]
\small
\begin{itemize}
  \item Mide estabilidad global de la proteína
  \item Se calcula respecto a estructura de referencia
  \item Plateau indica equilibrio
  \item Unidades: Å o nm
\end{itemize}
\end{tcolorbox}

\column{0.48\textwidth}
\begin{tcolorbox}[formulabox, title={RMSF (Fluctuación Cuadrática Media)}, equal height group=ehg-rmsd]
\[
\text{RMSF}_i = \sqrt{\langle \left|\mathbf{r}_i(t) - \langle\mathbf{r}_i\rangle\right|^2 \rangle_t}
\]
\small
\begin{itemize}
  \item Mide flexibilidad por residuo
  \item Regiones de alta RMSF = loops o terminales
  \item Correlaciona con factores B cristalográficos
  \item Herramienta: \texttt{gmx rmsf}, MDAnalysis
\end{itemize}
\end{tcolorbox}
\end{columns}
\end{frame}

%------------------------------------------------------------------------------
\begin{frame}{Otras Métricas Importantes}
\begin{columns}[T]
\column{0.48\textwidth}
\begin{tcolorbox}[formulabox, title={Radio de Giro ($R_g$)}, equal height group=ehg-met-a]
\small
\[
R_g(t) = \sqrt{\frac{\sum_i m_i |\mathbf{r}_i - \mathbf{r}_{cm}|^2}{\sum_i m_i}}
\]
Indica compacidad de la proteína. Una proteína plegada tiene $R_g$ estable.
\end{tcolorbox}
\vspace{0.4em}
\begin{tcolorbox}[definitionbox, title={Función de Distribución Radial}, equal height group=ehg-met-b]
\small
\[
g(r) = \frac{V}{N^2}\left\langle \sum_i \sum_{j\neq i} \delta(r - r_{ij}) \right\rangle
\]
Estructura de líquidos y solvatación.
\end{tcolorbox}

\column{0.48\textwidth}
\begin{tcolorbox}[warningbox, title={Puentes de Hidrógeno}, equal height group=ehg-met-a]
\small
Criterio geométrico:
\begin{itemize}
  \item $d_{D\cdots A} \leq 3.5$ Å
  \item $\angle_{D-H\cdots A} \geq 120°$
\end{itemize}
Indica estructura secundaria y estabilidad.
\end{tcolorbox}
\vspace{0.4em}
\begin{tcolorbox}[formulabox, title={Coeficiente de difusión}, equal height group=ehg-met-b]
\small
Einstein:
\[
D = \lim_{t\to\infty} \frac{\langle |\mathbf{r}(t) - \mathbf{r}(0)|^2 \rangle}{6t}
\]
\end{tcolorbox}
\end{columns}
\end{frame}

%==============================================================================
\section{Comparación de Programas}
%==============================================================================

\begin{frame}{Principales Programas de MD}

\begin{table}[h]
\centering
\scriptsize
\setlength{\tabcolsep}{5pt}
\renewcommand{\arraystretch}{1.25}
\begin{tabular}{@{}lcccc@{}}
\toprule
\textbf{Característica} & 
\cellcolor{PastelBlue!50}\textbf{GROMACS} & 
\cellcolor{PastelGreen!50}\textbf{OpenMM} & 
\cellcolor{PastelOrange!50}\textbf{NAMD} & 
\cellcolor{PastelPurple!40}\textbf{AMBER} \\
\midrule
Licencia        & LGPL (libre)      & MIT (libre)        & Libre (académico) & Comercial \\
Lenguaje        & C++/CUDA          & Python/C++/CUDA    & C++/TCL           & C/CUDA \\
GPU support     & \checkmark Excelente & \checkmark Excelente & \checkmark Bueno  & \checkmark Excelente \\
Interfaz        & Línea de comandos & Python API         & TCL/NAMD scripts  & Ambas \\
Velocidad       & Muy alta          & Alta               & Alta              & Muy alta \\
Escalabilidad   & Excelente HPC     & Media              & Excelente HPC     & Buena \\
Campos de F.    & AMBER/CHARMM/OPLS & Todos              & CHARMM/AMBER      & AMBER/GAFF \\
Análisis        & Herramientas propias & MDTraj/MDAnalysis & VMD/NAMD          & CPPTRAJ \\
Documentación   & Excelente         & Buena              & Buena             & Buena \\
Curva aprendizaje & Media           & Baja (Python)      & Media             & Media-alta \\
\midrule
\textbf{Mejor para} & Proteínas/membranas & Prototyping/ML & Simulaciones grandes & Bioquímica/AMBER FF \\
\bottomrule
\end{tabular}
\end{table}

\vspace{0.2em}
\begin{tcolorbox}[comparisonbox]
\scriptsize\centering
\textbf{Recomendación para principiantes:} GROMACS (documentación + tutoriales) o OpenMM (Python, mayor flexibilidad)
\end{tcolorbox}
\end{frame}

%------------------------------------------------------------------------------
\begin{frame}{Ecosistema de Herramientas MD}

\begin{center}
\begin{tikzpicture}[
  tool/.style={rectangle, rounded corners=4pt, minimum width=2.2cm, minimum height=0.55cm,
    text centered, font=\scriptsize, draw=DarkBlue!60},
  category/.style={rectangle, rounded corners=6pt, minimum width=2.4cm, minimum height=0.6cm,
    text centered, font=\small\bfseries, draw=DarkBlue, fill=DarkBlue, text=white},
  arrow/.style={->, DarkBlue!60}
]
% Categorías
\node[category] (prep) at (0,3) {Preparación};
\node[category] (sim) at (3.5,3) {Simulación};
\node[category] (anal) at (7,3) {Análisis};
\node[category] (vis) at (10.5,3) {Visualización};

% Herramientas
\node[tool, fill=PastelGreen!50] (pdbfixer) at (0,2) {PDBFixer};
\node[tool, fill=PastelGreen!40] (tleap) at (0,1.3) {tLEaP (AMBER)};
\node[tool, fill=PastelGreen!40] (pdb2gmx) at (0,0.6) {pdb2gmx (GROMACS)};
\node[tool, fill=PastelGreen!30] (charmm) at (0,-0.1) {CHARMM-GUI};

\node[tool, fill=PastelBlue!50] (gromacs) at (3.5,2) {GROMACS};
\node[tool, fill=PastelBlue!40] (openmm) at (3.5,1.3) {OpenMM};
\node[tool, fill=PastelBlue!40] (namd) at (3.5,0.6) {NAMD};
\node[tool, fill=PastelBlue!30] (amber) at (3.5,-0.1) {AMBER};

\node[tool, fill=PastelOrange!50] (mdanalysis) at (7,2) {MDAnalysis};
\node[tool, fill=PastelOrange!40] (mdtraj) at (7,1.3) {MDTraj};
\node[tool, fill=PastelOrange!40] (cpptraj) at (7,0.6) {CPPTRAJ};
\node[tool, fill=PastelOrange!30] (pytraj) at (7,-0.1) {pytraj};

\node[tool, fill=PastelPurple!50] (vmd) at (10.5,2) {VMD};
\node[tool, fill=PastelPurple!40] (pymol) at (10.5,1.3) {PyMOL};
\node[tool, fill=PastelPurple!40] (nglview) at (10.5,0.6) {NGLView};
\node[tool, fill=PastelPurple!30] (chimera) at (10.5,-0.1) {ChimeraX};

% Flechas principales
\draw[arrow, thick] (1.4,1) -- (2.3,1);
\draw[arrow, thick] (4.9,1) -- (5.8,1);
\draw[arrow, thick] (8.3,1) -- (9.3,1);

\end{tikzpicture}
\end{center}
\end{frame}

%==============================================================================
\section{Protocolo de Simulación de Proteínas}
%==============================================================================

\begin{frame}{Protocolo Completo: Simulación de Proteína}

\begin{columns}[T]
\column{0.5\textwidth}
\begin{enumerate}\small
  \item[\faCheckCircle] \textbf{Descarga}: PDB / AlphaFold2
  \item[\faCheckCircle] \textbf{Preparación}: PDBFixer, Modeller
  \item[\faCheckCircle] \textbf{Campo de fuerzas}: AMBER ff14SB / CHARMM36
  \item[\faCheckCircle] \textbf{Caja}: dodecaedro rómbico, 1.0 nm margen
  \item[\faCheckCircle] \textbf{Solvatar}: SPC/E o TIP3P
  \item[\faCheckCircle] \textbf{Neutralizar}: Na$^+$/Cl$^-$ a 0.15 M
  \item[\faCheckCircle] \textbf{Minimizar}: 5000 pasos steepest descent
  \item[\faCheckCircle] \textbf{Equilibrar NVT}: 100 ps, $T = 300$ K
  \item[\faCheckCircle] \textbf{Equilibrar NPT}: 100 ps, $P = 1$ bar
  \item[\faCheckCircle] \textbf{Producción}: 100 ns--1 $\mu$s
\end{enumerate}

\column{0.47\textwidth}
\begin{tcolorbox}[formulabox, title={Parámetros típicos GROMACS}]
\scriptsize\ttfamily
integrator = md \\
nsteps = 50000000 ; 100 ns \\
dt = 0.002 ; 2 fs \\
\textcolor{DarkGreen}{; Temperatura} \\
tcoupl = V-rescale \\
ref\_t = 300 \\
tau\_t = 0.1 \\
\textcolor{DarkGreen}{; Presión} \\
pcoupl = Parrinello-Rahman \\
ref\_p = 1.0 \\
tau\_p = 2.0 \\
\textcolor{DarkGreen}{; Interacciones} \\
rcoulomb = 1.2 \\
rvdw = 1.2 \\
coulombtype = PME
\end{tcolorbox}
\end{columns}
\end{frame}

%==============================================================================
\section{Técnicas Avanzadas}
%==============================================================================

\begin{frame}{Técnicas Avanzadas de MD}
\begin{columns}[T]
\column{0.48\textwidth}
\begin{block}{Muestreo mejorado}
  \begin{itemize}
    \item \textbf{Metadinámica}: potencial de inundación gaussiana para superar barreras energéticas
    \item \textbf{Umbrella sampling}: histogramas de fuerza con coordenadas de reacción
    \item \textbf{Replica Exchange MD}: múltiples temperaturas en paralelo
    \item \textbf{Steered MD}: fuerza externa para estudiar desenrollamiento
  \end{itemize}
\end{block}

\column{0.48\textwidth}
\begin{block}{Cálculo de propiedades termodinámicas}
  \begin{itemize}
    \item \textbf{Perturbación de energía libre (FEP)}: $\Delta G = -k_BT\ln\langle e^{-\beta\Delta U}\rangle$
    \item \textbf{Integración termodinámica (TI)}: $\Delta G = \int_0^1 \langle \partial H/\partial\lambda\rangle\,d\lambda$
    \item \textbf{MM-PBSA/GBSA}: energía libre de unión
    \item \textbf{PMF}: potencial de fuerza media
  \end{itemize}
\end{block}
\end{columns}
\vspace{0.4em}
\begin{tcolorbox}[formulabox, title={Energía libre de Gibbs (FEP)}]
\[
\Delta G_{A\to B} = -k_B T \ln \left\langle \exp\!\left(-\frac{\Delta U_{A\to B}}{k_B T}\right)\right\rangle_A
\]
\end{tcolorbox}
\end{frame}

%------------------------------------------------------------------------------
\begin{frame}{MD con Aprendizaje Automático}
\begin{columns}[T]
\column{0.52\textwidth}
\begin{block}{Machine Learning en MD}
  \begin{itemize}
    \item \textbf{Campos de fuerzas ML}: NequIP, MACE, ANI — precisión cuántica con velocidad clásica
    \item \textbf{Potenciales neurales}: redes grafos moleculares para predicción de fuerzas
    \item \textbf{MSM} (Markov State Models): identificación de estados conformacionales
    \item \textbf{OpenMM-ML}: integración directa con PyTorch/TensorFlow
  \end{itemize}
\end{block}

\column{0.45\textwidth}
\begin{tcolorbox}[warningbox, title={Flujo ML-MD}]
\small
\begin{enumerate}
  \item Generar datos QM (DFT/CCSD)
  \item Entrenar potencial neural
  \item Validar contra benchmarks
  \item Simulación MD con potencial ML
  \item Análisis y propiedades
\end{enumerate}
\end{tcolorbox}
\end{columns}
\end{frame}

%==============================================================================
\section{Resumen y Perspectivas}
%==============================================================================

\begin{frame}[shrink=3]{Resumen: Lo que Aprendimos}

\begin{columns}[T]
\column{0.48\textwidth}
\begin{tcolorbox}[definitionbox, title={\faAtom\ Fundamentos}, equal height group=ehg-res-a]
\small
\begin{itemize}
  \item[$\checkmark$] Ecuaciones de Newton y campo de fuerzas
  \item[$\checkmark$] Potencial de Lennard-Jones
  \item[$\checkmark$] Distribución de Maxwell-Boltzmann
  \item[$\checkmark$] Integradores: Verlet, V-Verlet, Leap-Frog
\end{itemize}
\end{tcolorbox}
\vspace{0.3em}
\begin{tcolorbox}[formulabox, title={\faFlask\ Práctica}, equal height group=ehg-res-b]
\small
\begin{itemize}
  \item[$\checkmark$] PBC y convención de imagen mínima
  \item[$\checkmark$] Ensambles NVE, NVT, NPT
  \item[$\checkmark$] Termostatos y barostatos
  \item[$\checkmark$] Preparación y protocolo de simulación
\end{itemize}
\end{tcolorbox}

\column{0.48\textwidth}
\begin{tcolorbox}[warningbox, title={\faChartBar\ Análisis}, equal height group=ehg-res-a]
\small
\begin{itemize}
  \item[$\checkmark$] RMSD, RMSF, $R_g$, RDF
  \item[$\checkmark$] Puentes de hidrógeno
  \item[$\checkmark$] Coeficiente de difusión
  \item[$\checkmark$] MDAnalysis y MDTraj (Python)
\end{itemize}
\end{tcolorbox}
\vspace{0.3em}
\begin{tcolorbox}[definitionbox, title={\faLaptopCode\ Herramientas}, equal height group=ehg-res-b]
\small
\begin{itemize}
  \item[$\checkmark$] GROMACS, OpenMM, NAMD, AMBER
  \item[$\checkmark$] Campos de fuerzas: AMBER, CHARMM, OPLS
  \item[$\checkmark$] Modelos de agua: SPC/E, TIP3P
  \item[$\checkmark$] Técnicas avanzadas: FEP, metadinámica
\end{itemize}
\end{tcolorbox}
\end{columns}
\end{frame}

%------------------------------------------------------------------------------
\begin{frame}{Perspectivas y Tendencias Actuales}
\begin{columns}[T]
\column{0.48\textwidth}
\begin{block}{Tendencias 2024--2025}
  \begin{itemize}
    \item \textbf{MD asistida por IA}: AlphaFold + MD
    \item \textbf{Potenciales neurales}: MACE, NequIP, ANI2x
    \item \textbf{Simulaciones de larga escala}: ms con Anton 3
    \item \textbf{Coarse-graining con ML}: modelos Martini + ML
    \item \textbf{Enhanced sampling}: OPES, on-the-fly ML
  \end{itemize}
\end{block}

\column{0.48\textwidth}
\begin{block}{Aplicaciones emergentes}
  \begin{itemize}
    \item Diseño racional de fármacos
    \item Dinámica de anticuerpos y nanoanticuerpos
    \item Materiales 2D y polímeros conductores
    \item Sistemas de membrana complejos
    \item Simulación de virus completos
  \end{itemize}
\end{block}
\end{columns}

\vspace{0.5em}
\begin{tcolorbox}[comparisonbox]
\centering\small
\textbf{Recursos recomendados:} GROMACS tutorials (Lemkul), OpenMM cookbook, \textit{Understanding Molecular Simulation} (Frenkel \& Smit), \textit{Computer Simulation of Liquids} (Allen \& Tildesley)
\end{tcolorbox}
\end{frame}

%==============================================================================
\section{Instalación de GROMACS}
%==============================================================================

%------------------------------------------------------------------------------
\begin{frame}[fragile,shrink=8]{Instalación de GROMACS en Ubuntu/Linux}
\begin{columns}[T]
\column{0.48\textwidth}
\begin{tcolorbox}[definitionbox, title={Prerrequisitos}]
\begin{lstlisting}[style=bashstyle]
sudo apt update && sudo apt install -y \
  build-essential cmake wget \
  libfftw3-dev libopenmpi-dev \
  openmpi-bin libhwloc-dev
# Verificar versiones minimas
cmake --version   # >= 3.18
gcc --version     # >= 7
\end{lstlisting}
\end{tcolorbox}
\vspace{0.3em}
\begin{tcolorbox}[warningbox, title={Opción rápida: paquete apt}]
\begin{lstlisting}[style=bashstyle]
# Ubuntu 22.04/24.04 (version del repo)
sudo apt install -y gromacs gromacs-data
gmx --version
# Puede ser una version mas antigua (2022.x)
\end{lstlisting}
\end{tcolorbox}
\column{0.48\textwidth}
\begin{tcolorbox}[formulabox, title={Opción recomendada: compilar GROMACS 2024.2}]
\begin{lstlisting}[style=bashstyle]
wget https://ftp.gromacs.org/gromacs/gromacs-2024.2.tar.gz
tar -xzf gromacs-2024.2.tar.gz
cd gromacs-2024.2 && mkdir build && cd build
cmake .. \
  -DGMX_BUILD_OWN_FFTW=ON \
  -DGMX_GPU=OFF \
  -DCMAKE_BUILD_TYPE=Release
make -j4
sudo make install
# Activar en sesion actual
source /usr/local/gromacs/bin/GMXRC
# Activar permanentemente
echo 'source /usr/local/gromacs/bin/GMXRC' \
  >> ~/.bashrc && source ~/.bashrc
gmx --version
\end{lstlisting}
\end{tcolorbox}
\begin{tcolorbox}[comparisonbox]
\small\textbf{Con GPU (CUDA):} añadir \texttt{-DGMX\_GPU=CUDA} \\y tener instalado \texttt{nvidia-cuda-toolkit}
\end{tcolorbox}
\end{columns}
\end{frame}

%------------------------------------------------------------------------------
\begin{frame}[fragile,shrink=8]{Instalación de GROMACS en Google Colab}
\begin{columns}[T]
\column{0.50\textwidth}
\begin{tcolorbox}[warningbox, title={Opción 1: conda-forge (recomendada, \textasciitilde{}2 min)}]
\begin{lstlisting}[style=bashstyle]
# Celda 1 (Python): instalar condacolab
!pip install -q condacolab
import condacolab
condacolab.install()  # reinicia el kernel
\end{lstlisting}
\begin{lstlisting}[style=bashstyle]
# Celda 2 (Python): instalar GROMACS
!conda install -c conda-forge gromacs -y
!gmx --version
\end{lstlisting}
\end{tcolorbox}
\vspace{0.3em}
\begin{tcolorbox}[definitionbox, title={Prerrequisitos (compilación directa)}]
\begin{lstlisting}[style=bashstyle]
%%bash
apt-get update -qq
apt-get install -y cmake build-essential \
  libfftw3-dev libopenmpi-dev \
  openmpi-bin wget -qq
\end{lstlisting}
\end{tcolorbox}
\column{0.47\textwidth}
\begin{tcolorbox}[formulabox, title={Opción 2: compilar en Colab (\textasciitilde{}10 min)}]
\begin{lstlisting}[style=bashstyle]
%%bash
wget -q https://ftp.gromacs.org/gromacs/gromacs-2024.2.tar.gz
tar -xzf gromacs-2024.2.tar.gz
cd gromacs-2024.2 && mkdir build && cd build
cmake .. \
  -DGMX_BUILD_OWN_FFTW=ON \
  -DGMX_GPU=OFF \
  -DCMAKE_INSTALL_PREFIX=/usr/local \
  -DCMAKE_BUILD_TYPE=Release
make -j4 -s && make install -s
\end{lstlisting}
\end{tcolorbox}
\begin{tcolorbox}[comparisonbox]
\begin{lstlisting}[style=bashstyle]
# Celda Python: integrar con subprocesos
import subprocess, os
os.environ['PATH'] += \
  ':/usr/local/gromacs/bin'
r = subprocess.run(
  ['gmx', '--version'],
  capture_output=True, text=True)
print(r.stdout.split('\n')[1])
\end{lstlisting}
\end{tcolorbox}
\end{columns}
\end{frame}

%==============================================================================
\section{Tutorial GROMACS: Dinámica de Lisozima}
%==============================================================================

%------------------------------------------------------------------------------
\begin{frame}{Tutorial: Simulación de Lisozima con GROMACS}
\begin{columns}[T]
\column{0.48\textwidth}
\begin{block}{Sistema}
  \begin{itemize}\small
    \item \textbf{Proteína}: Lisozima (1AKI, 129 residuos)
    \item \textbf{Campo de fuerzas}: AMBER99SB-ILDN
    \item \textbf{Solvente}: TIP3P, caja dodecaédrica
    \item \textbf{Iones}: Na$^+$/Cl$^-$ (neutralización + 0.15\,M)
    \item \textbf{Temperatura}: 300\,K \quad \textbf{Presión}: 1\,bar
    \item \textbf{Producción}: 1\,ns ($\Delta t = 2$\,fs, 500\,000 pasos)
  \end{itemize}
\end{block}
\column{0.48\textwidth}
\begin{block}{Protocolo de 8 pasos}
  \begin{enumerate}\small
    \item[\textbf{1.}] \textbf{pdb2gmx}: topología + parámetros FF
    \item[\textbf{2.}] \textbf{editconf + solvate}: caja + agua
    \item[\textbf{3.}] \textbf{genion}: iones + neutralización
    \item[\textbf{4.}] \textbf{EM}: minimización (steepest descent)
    \item[\textbf{5.}] \textbf{NVT eq.}: 100\,ps a $T$ constante
    \item[\textbf{6.}] \textbf{NPT eq.}: 100\,ps a $T$, $P$ constantes
    \item[\textbf{7.}] \textbf{MD}: producción 1\,ns
    \item[\textbf{8.}] \textbf{Análisis}: RMSD, RMSF, $R_g$\ldots
  \end{enumerate}
\end{block}
\end{columns}
\vspace{0.3em}
\begin{tcolorbox}[comparisonbox]
  \centering\small
  Ref: Lemkul JA (2019) \textit{J.\,Phys.\,Chem.\,B} 123:6254 \quad
  \texttt{tutorials.gromacs.org} \quad GROMACS 2023
\end{tcolorbox}
\end{frame}

%------------------------------------------------------------------------------
\begin{frame}[fragile]{Pasos 1--2: Preparación del PDB y Solvatación}
\begin{columns}[T]
\column{0.50\textwidth}
\begin{tcolorbox}[definitionbox, title={1. Obtener y limpiar PDB}]
\begin{lstlisting}[style=bashstyle]
# Descargar lisozima de RCSB
wget https://files.rcsb.org/download/1AKI.pdb
# Eliminar agua cristalografica
grep -v HOH 1AKI.pdb > 1AKI_clean.pdb
\end{lstlisting}
\end{tcolorbox}
\begin{tcolorbox}[definitionbox, title={2. Generar topologia (FF + agua)}]
\begin{lstlisting}[style=bashstyle]
gmx pdb2gmx -f 1AKI_clean.pdb \
  -o protein.gro -water tip3p \
  -ff amber99sb-ildn
# Genera: topol.top, posre.itp, protein.gro
\end{lstlisting}
\end{tcolorbox}
\column{0.47\textwidth}
\begin{tcolorbox}[definitionbox, title={3. Caja dodecaedrica (1 nm margen)}]
\begin{lstlisting}[style=bashstyle]
gmx editconf -f protein.gro \
  -o boxed.gro -c -d 1.0 \
  -bt dodecahedron
\end{lstlisting}
\end{tcolorbox}
\begin{tcolorbox}[definitionbox, title={4. Solvatacion con TIP3P}]
\begin{lstlisting}[style=bashstyle]
gmx solvate -cp boxed.gro \
  -cs spc216.gro \
  -o solvated.gro -p topol.top
# Actualiza topol.top: agrega SOL
\end{lstlisting}
\end{tcolorbox}
\end{columns}
\end{frame}

%------------------------------------------------------------------------------
\begin{frame}[fragile,shrink=2]{Paso 3--4: Iones y Minimizacion de Energia}
\begin{columns}[T]
\column{0.50\textwidth}
\begin{tcolorbox}[definitionbox, title={3. Anadir iones (neutralizar)}]
\begin{lstlisting}[style=bashstyle]
# ions.mdp: integrator=steep, nsteps=0
gmx grompp -f ions.mdp \
  -c solvated.gro -p topol.top \
  -o ions.tpr -maxwarn 1
echo "SOL" | gmx genion \
  -s ions.tpr -o system.gro \
  -p topol.top -pname NA \
  -nname CL -neutral -conc 0.15
\end{lstlisting}
\end{tcolorbox}
\column{0.47\textwidth}
\begin{tcolorbox}[warningbox, title={4. minim.mdp (puntos clave)}]
\begin{lstlisting}[style=mdpstyle]
integrator   = steep
emtol        = 1000.0   ; kJ/mol/nm
emstep       = 0.01     ; nm
nsteps       = 50000
nstlist      = 1
cutoff-scheme = Verlet
coulombtype  = PME
rcoulomb     = 1.0
rvdw         = 1.0
\end{lstlisting}
\end{tcolorbox}
\begin{tcolorbox}[definitionbox, title={Ejecutar EM}]
\begin{lstlisting}[style=bashstyle]
gmx grompp -f minim.mdp \
  -c system.gro -p topol.top -o em.tpr
gmx mdrun -v -deffnm em
# Verificar: Epot debe ser negativo
echo "Potential" | gmx energy \
  -f em.edr -o potential.xvg
\end{lstlisting}
\end{tcolorbox}
\end{columns}
\end{frame}

%------------------------------------------------------------------------------
\begin{frame}[fragile,shrink=3]{Paso 5: Equilibracion NVT (100 ps)}
\begin{columns}[T]
\column{0.51\textwidth}
\begin{tcolorbox}[warningbox, title={nvt.mdp --- parametros clave}]
\begin{lstlisting}[style=mdpstyle]
integrator  = md
dt          = 0.002      ; 2 fs
nsteps      = 50000      ; 100 ps
nstxout-compressed = 500 ; cada 1 ps
; Restricciones de enlaces H
constraints         = h-bonds
constraint-algorithm = lincs
; Termostato V-rescale
tcoupl    = V-rescale
tc-grps   = Protein Non-Protein
tau-t     = 0.1  0.1     ; ps
ref-t     = 300  300     ; K
; Electrostatos
coulombtype  = PME
rcoulomb     = 1.0
rvdw         = 1.0
DispCorr     = EnerPres
; Restricciones de posicion
define    = -DPOSRES
; Velocidades iniciales
gen-vel   = yes
gen-temp  = 300
\end{lstlisting}
\end{tcolorbox}
\column{0.46\textwidth}
\begin{tcolorbox}[definitionbox, title={Ejecutar NVT}]
\begin{lstlisting}[style=bashstyle]
gmx grompp -f nvt.mdp \
  -c em.gro -r em.gro \
  -p topol.top -o nvt.tpr
gmx mdrun -v -deffnm nvt
\end{lstlisting}
\end{tcolorbox}
\begin{tcolorbox}[definitionbox, title={Verificar temperatura}]
\begin{lstlisting}[style=bashstyle]
echo "Temperature" | gmx energy \
  -f nvt.edr -o temperature.xvg
xmgrace temperature.xvg
# Debe estabilizarse en ~300 K
\end{lstlisting}
\end{tcolorbox}
\begin{tcolorbox}[comparisonbox]
\small\textbf{Nota}: \texttt{-DPOSRES} mantiene la
proteina fija durante la equilibracion
termica para evitar distorsiones
estructurales iniciales.
\end{tcolorbox}
\end{columns}
\end{frame}

%------------------------------------------------------------------------------
\begin{frame}[fragile,shrink=3]{Pasos 6--7: Equilibracion NPT y Produccion}
\begin{columns}[T]
\column{0.50\textwidth}
\begin{tcolorbox}[warningbox, title={npt.mdp (cambios vs nvt.mdp)}]
\begin{lstlisting}[style=mdpstyle]
; Barostato Parrinello-Rahman
pcoupl         = Parrinello-Rahman
pcoupltype     = isotropic
tau-p          = 2.0     ; ps
ref-p          = 1.0     ; bar
compressibility = 4.5e-5 ; bar^-1
; NO generar velocidades nuevas
gen-vel   = no
\end{lstlisting}
\end{tcolorbox}
\begin{lstlisting}[style=bashstyle]
gmx grompp -f npt.mdp -c nvt.gro \
  -r nvt.gro -t nvt.cpt \
  -p topol.top -o npt.tpr
gmx mdrun -v -deffnm npt
echo "Pressure" | gmx energy \
  -f npt.edr -o pressure.xvg
echo "Density" | gmx energy \
  -f npt.edr -o density.xvg
\end{lstlisting}
\column{0.47\textwidth}
\begin{tcolorbox}[formulabox, title={md.mdp --- produccion 1 ns}]
\begin{lstlisting}[style=mdpstyle]
integrator  = md
dt          = 0.002
nsteps      = 500000     ; 1 ns
; Sin restricciones de posicion
; (eliminar -DPOSRES)
nstxout-compressed = 500
nstlog      = 500
nstenergy   = 500
; Continuar desde NPT
gen-vel   = no
\end{lstlisting}
\end{tcolorbox}
\begin{lstlisting}[style=bashstyle]
gmx grompp -f md.mdp -c npt.gro \
  -t npt.cpt -p topol.top -o md.tpr
# Produccion (~4 h en 4 cores)
gmx mdrun -v -deffnm md \
  -ntmpi 1 -ntomp 4
\end{lstlisting}
\end{columns}
\end{frame}

%------------------------------------------------------------------------------
\begin{frame}[fragile]{Analisis I: RMSD, RMSF y Energia}
\begin{columns}[T]
\column{0.50\textwidth}
\begin{tcolorbox}[definitionbox, title={RMSD del backbone}]
\begin{lstlisting}[style=bashstyle]
# Grupo 4 = backbone (C, CA, N)
echo "4 4" | gmx rms \
  -s md.tpr -f md.xtc \
  -o rmsd.xvg -tu ns
xmgrace rmsd.xvg
# Valor esperado: 0.1-0.3 nm
\end{lstlisting}
\end{tcolorbox}
\begin{tcolorbox}[definitionbox, title={RMSF por residuo}]
\begin{lstlisting}[style=bashstyle]
echo "4" | gmx rmsf \
  -s md.tpr -f md.xtc \
  -o rmsf.xvg -res
# Picos = loops flexibles
\end{lstlisting}
\end{tcolorbox}
\column{0.47\textwidth}
\begin{tcolorbox}[definitionbox, title={Energia, temperatura, presion}]
\begin{lstlisting}[style=bashstyle]
echo "Potential" | gmx energy \
  -f md.edr -o energy.xvg
echo "Temperature" | gmx energy \
  -f md.edr -o temp.xvg
echo "Pressure" | gmx energy \
  -f md.edr -o pres.xvg
echo "Density" | gmx energy \
  -f md.edr -o density.xvg
\end{lstlisting}
\end{tcolorbox}
\begin{tcolorbox}[comparisonbox]
\small
\textbf{Valores de referencia}\\
\textbf{RMSD}: 0.1--0.3\,nm (estable)\\
\textbf{T}: $300 \pm 5$\,K\\
\textbf{P}: $1 \pm 100$\,bar\\
\textbf{$\rho$}: $\approx 1000$\,kg/m$^3$
\end{tcolorbox}
\end{columns}
\end{frame}

%------------------------------------------------------------------------------
\begin{frame}[fragile,shrink=16]{Analisis II: $R_g$, H-bonds, SASA y Clustering}
\begin{columns}[T]
\column{0.50\textwidth}
\begin{tcolorbox}[definitionbox, title={Radio de giro ($R_g$)}]
\begin{lstlisting}[style=bashstyle]
echo "1" | gmx gyrate \
  -s md.tpr -f md.xtc \
  -o gyrate.xvg
# Rg constante = estructura compacta
\end{lstlisting}
\end{tcolorbox}
\begin{tcolorbox}[definitionbox, title={Puentes de hidrogeno}]
\begin{lstlisting}[style=bashstyle]
# H-bonds intraproteina
echo "1 1" | gmx hbond \
  -s md.tpr -f md.xtc \
  -num hbnum_prot.xvg
# H-bonds proteina-agua
echo "1 12" | gmx hbond \
  -s md.tpr -f md.xtc \
  -num hbnum_wat.xvg
\end{lstlisting}
\end{tcolorbox}
\column{0.47\textwidth}
\begin{tcolorbox}[definitionbox, title={SASA (superficie accesible)}]
\begin{lstlisting}[style=bashstyle]
echo "1" | gmx sasa \
  -s md.tpr -f md.xtc \
  -o sasa.xvg -probe 0.14
# probe = radio agua (0.14 nm)
\end{lstlisting}
\end{tcolorbox}
\begin{tcolorbox}[definitionbox, title={Clustering de conformaciones}]
\begin{lstlisting}[style=bashstyle]
echo "4 4" | gmx cluster \
  -s md.tpr -f md.xtc \
  -method gromos -cutoff 0.2 \
  -cl clusters.pdb \
  -clndx cluster.ndx
\end{lstlisting}
\end{tcolorbox}
\begin{tcolorbox}[warningbox, title={Visualizacion con VMD/PyMOL}]
\begin{lstlisting}[style=bashstyle]
gmx trjconv -f md.xtc -s md.tpr \
  -o traj_vis.pdb -dt 100 -pbc mol
\end{lstlisting}
\end{tcolorbox}
\end{columns}
\end{frame}

%==============================================================================
\section{Tutorial 2: Complejo Proteína-Ligando}
%==============================================================================

%------------------------------------------------------------------------------
\begin{frame}{Tutorial 2: MD de un Complejo Proteína-Ligando}
\begin{columns}[T]
\column{0.48\textwidth}
\begin{block}{Sistema}
  \begin{itemize}\small
    \item \textbf{Proteína}: T4 Lisozima L99A (PDB: 3HTB, 162 res.)
    \item \textbf{Ligando}: JZ4 (benzimidazol-2-ona, \ch{C7H6N2O})
    \item \textbf{Campo de fuerzas}: AMBER99SB-ILDN + GAFF2
    \item \textbf{Solvente}: TIP3P, caja cúbica
    \item \textbf{Iones}: Na$^+$/Cl$^-$ a 0.15\,M
    \item \textbf{Producción}: 100\,ns ($\Delta t = 2$\,fs)
  \end{itemize}
\end{block}
\column{0.48\textwidth}
\begin{block}{Flujo de trabajo}
  \begin{enumerate}\small
    \item[\textbf{1.}] Preparar PDB y separar ligando
    \item[\textbf{2.}] Parametrizar ligando (ACPYPE/GAFF2)
    \item[\textbf{3.}] Combinar topologías proteína + ligando
    \item[\textbf{4.}] Caja, solvatación e iones
    \item[\textbf{5.}] EM, equilibración NVT y NPT
    \item[\textbf{6.}] Producción MD
    \item[\textbf{7.}] Análisis: RMSD, contactos, MM-PBSA
  \end{enumerate}
\end{block}
\end{columns}
\vspace{0.3em}
\begin{tcolorbox}[comparisonbox]
\centering\small
Ref: Lemkul JA (2014) \textit{GROMACS Tutorial: Protein--Ligand Complex} \quad
Ligando JZ4 = sitio hidrofóbico de T4 lisozima L99A \quad \texttt{www.mdtutorials.com}
\end{tcolorbox}
\end{frame}

%------------------------------------------------------------------------------
\begin{frame}[fragile,shrink=5]{Pasos 1--2: Preparación y Parametrización del Ligando}
\begin{columns}[T]
\column{0.48\textwidth}
\begin{tcolorbox}[definitionbox, title={1. Obtener PDB y separar ligando}]
\begin{lstlisting}[style=bashstyle]
# Descargar complejo cristalografico
wget https://files.rcsb.org/download/3HTB.pdb
# Extraer proteina (solo ATOM)
grep "^ATOM" 3HTB.pdb > protein.pdb
# Extraer ligando JZ4
grep "^HETATM" 3HTB.pdb | \
  grep " JZ4 " > jz4.pdb
\end{lstlisting}
\end{tcolorbox}
\vspace{0.3em}
\begin{tcolorbox}[warningbox, title={2. Instalar ACPYPE y AmberTools}]
\begin{lstlisting}[style=bashstyle]
# Via conda-forge (recomendado)
conda install -c conda-forge \
  acpype ambertools openbabel -y
# Via pip
pip install acpype
# Verificar
acpype --version
\end{lstlisting}
\end{tcolorbox}
\column{0.48\textwidth}
\begin{tcolorbox}[formulabox, title={3. Generar parámetros GAFF2 (ACPYPE)}]
\begin{lstlisting}[style=bashstyle]
# Convertir PDB a MOL2 (asigna H)
obabel jz4.pdb -O jz4.mol2 -h
# Parametrizar con cargas AM1-BCC
# y campo de fuerzas GAFF2
acpype -i jz4.mol2 -c bcc -a gaff2 -n 0
# -n 0 = carga neta del ligando
# Archivos generados en jz4.acpype/:
#   jz4_GMX.itp  (parametros GROMACS)
#   jz4_GMX.gro  (coordenadas)
#   jz4_GMX.top  (topologia)
ls jz4.acpype/
\end{lstlisting}
\end{tcolorbox}
\begin{tcolorbox}[comparisonbox]
\small\textbf{Cargas AM1-BCC}: rápidas y adecuadas para MD.\\
\textbf{Mayor precisión}: RESP con GAUSSIAN/ORCA + Multiwfn.\\
\textbf{Alternativa}: CGenFF (CHARMM-GUI) para ff CHARMM.
\end{tcolorbox}
\end{columns}
\end{frame}

%------------------------------------------------------------------------------
\begin{frame}[fragile,shrink=5]{Paso 3: Topología del Complejo}
\begin{columns}[T]
\column{0.50\textwidth}
\begin{tcolorbox}[definitionbox, title={3a. Topología de la proteína}]
\begin{lstlisting}[style=bashstyle]
gmx pdb2gmx -f protein.pdb \
  -o protein.gro -water tip3p \
  -ff amber99sb-ildn -ignh
# Genera: topol.top, posre.itp, protein.gro
\end{lstlisting}
\end{tcolorbox}
\vspace{0.3em}
\begin{tcolorbox}[warningbox, title={3c. Editar topol.top (añadir ligando)}]
\begin{lstlisting}[style=mdpstyle]
; Agregar antes de [ system ]
#include "jz4.acpype/jz4_GMX.itp"
; Posicion restraints del ligando
#ifdef POSRES_LIG
#include "posre_jz4.itp"
#endif
; Actualizar [ molecules ]
Protein_chain_A   1
JZ4               1
\end{lstlisting}
\end{tcolorbox}
\column{0.47\textwidth}
\begin{tcolorbox}[formulabox, title={3b. Combinar coordenadas}]
\begin{lstlisting}[style=bashstyle]
# Copiar .gro del ligando
cp jz4.acpype/jz4_GMX.gro jz4.gro
# Combinar proteina + ligando
# (ajustar natoms en la cabecera)
python3 << 'EOF'
with open('protein.gro') as f: p=f.readlines()
with open('jz4.gro') as f: l=f.readlines()
n = int(p[1]) + int(l[1])
with open('complex.gro','w') as out:
  out.write('Complex protlig\n')
  out.write(f'{n}\n')
  out.writelines(p[2:-1])
  out.writelines(l[2:-1])
  out.write(p[-1])
EOF
\end{lstlisting}
\end{tcolorbox}
\begin{tcolorbox}[comparisonbox]
\small\textbf{Orden en topol.top}: el \texttt{\#include} del ligando
debe ir \textbf{antes} de \texttt{[ system ]}. El orden en \texttt{[ molecules ]}
debe coincidir exactamente con \texttt{complex.gro}.
\end{tcolorbox}
\end{columns}
\end{frame}

%------------------------------------------------------------------------------
\begin{frame}[fragile,shrink=5]{Pasos 4--5: Sistema, Minimización y Equilibración}
\begin{columns}[T]
\column{0.50\textwidth}
\begin{tcolorbox}[definitionbox, title={4. Caja, solvatación e iones}]
\begin{lstlisting}[style=bashstyle]
# Caja cubica, margen 1.2 nm
gmx editconf -f complex.gro \
  -o boxed.gro -c -d 1.2 -bt cubic
# Solvatar con TIP3P
gmx solvate -cp boxed.gro \
  -cs spc216.gro \
  -o solvated.gro -p topol.top
# Agregar iones (0.15 M NaCl)
gmx grompp -f ions.mdp \
  -c solvated.gro -p topol.top \
  -o ions.tpr -maxwarn 2
echo "SOL" | gmx genion \
  -s ions.tpr -o system.gro \
  -p topol.top -pname NA \
  -nname CL -neutral -conc 0.15
\end{lstlisting}
\end{tcolorbox}
\column{0.47\textwidth}
\begin{tcolorbox}[formulabox, title={5. EM, NVT y NPT}]
\begin{lstlisting}[style=bashstyle]
# Minimizacion de energia
gmx grompp -f minim.mdp \
  -c system.gro -p topol.top -o em.tpr
gmx mdrun -v -deffnm em
# Equilibracion NVT 100 ps
gmx grompp -f nvt.mdp \
  -c em.gro -r em.gro \
  -p topol.top -o nvt.tpr
gmx mdrun -v -deffnm nvt
# Equilibracion NPT 100 ps
gmx grompp -f npt.mdp \
  -c nvt.gro -r nvt.gro \
  -t nvt.cpt -p topol.top -o npt.tpr
gmx mdrun -v -deffnm npt
\end{lstlisting}
\end{tcolorbox}
\begin{tcolorbox}[warningbox, title={Restricciones del ligando}]
\begin{lstlisting}[style=bashstyle]
# Generar posre para el ligando
gmx genrestr -f jz4_GMX.gro \
  -o posre_jz4.itp \
  -fc 1000 1000 1000
\end{lstlisting}
\end{tcolorbox}
\end{columns}
\end{frame}

%------------------------------------------------------------------------------
\begin{frame}[fragile,shrink=5]{Análisis del Complejo Proteína-Ligando}
\begin{columns}[T]
\column{0.50\textwidth}
\begin{tcolorbox}[definitionbox, title={RMSD proteína y ligando}]
\begin{lstlisting}[style=bashstyle]
# Crear grupo para el ligando
gmx make_ndx -f md.tpr -o index.ndx
# Dentro de make_ndx:
# > r JZ4    (selecciona residuo JZ4)
# > name X JZ4_lig
# > q
# RMSD backbone proteina
echo "4 4" | gmx rms -s md.tpr \
  -f md.xtc -o rmsd_prot.xvg -tu ns
# RMSD ligando
echo "JZ4_lig JZ4_lig" | gmx rms \
  -s md.tpr -f md.xtc \
  -n index.ndx -o rmsd_lig.xvg -tu ns
\end{lstlisting}
\end{tcolorbox}
\begin{tcolorbox}[definitionbox, title={Distancia y contactos}]
\begin{lstlisting}[style=bashstyle]
# Distancia minima proteina-ligando
echo "1 JZ4_lig" | gmx mindist \
  -s md.tpr -f md.xtc \
  -n index.ndx -od mindist.xvg
\end{lstlisting}
\end{tcolorbox}
\column{0.47\textwidth}
\begin{tcolorbox}[warningbox, title={Energía libre: MM-PBSA}]
\begin{lstlisting}[style=bashstyle]
# Instalar gmx_MMPBSA
conda install -c conda-forge \
  gmx_mmpbsa ambertools -y
# input file mmpbsa.in
# &general startframe=1, endframe=500,
#   interval=5, verbose=2, /
# Ejecutar MM-PBSA
gmx_MMPBSA -O -i mmpbsa.in \
  -cs md.tpr -ct md.xtc \
  -ci index.ndx -cg 1 JZ4_lig \
  -cp topol.top \
  -o FINAL_RESULTS.dat
# Visualizar resultados
gmx_MMPBSA_ana -f FINAL_RESULTS.dat
\end{lstlisting}
\end{tcolorbox}
\begin{tcolorbox}[comparisonbox]
\small
$\Delta G_{bind}$ negativo $\rightarrow$ unión estable.\\Valores típicos para inhibidores:\\$\Delta G_{bind} \approx -20$ a $-60$\,kJ/mol.\\\textbf{Descomposición por residuos} identifica\\los hotspots del sitio de unión.
\end{tcolorbox}
\end{columns}
\end{frame}

%------------------------------------------------------------------------------
\begin{frame}[plain]
\begin{center}
{\Large\bfseries\color{DarkBlue} ¡Gracias!}\\[0.8em]
{\large\color{SubtitleColor} Dinámica Molecular: Principios, Métodos y Aplicaciones}\\[1.2em]

\begin{tikzpicture}[scale=1.0]
\foreach \x/\y/\c in {
  0/0/PastelBlue, 1.5/0.4/PastelGreen, 0.7/1.2/PastelPurple,
  2.5/0.9/PastelOrange, 2.0/1.8/PastelTeal, 3.5/0.2/PastelRed,
  3.8/1.4/PastelYellow, 0.3/2.2/PastelBlue!60, 4.5/0.8/PastelGreen!60}{
  \node[circle, fill=\c, draw=DarkBlue!40, minimum size=16pt] at (\x,\y) {};
}
\draw[->, thick, DarkBlue!70] (0,0) -- (0.5, 0.4);
\draw[->, thick, DarkGreen!70] (1.5,0.4) -- (1.0, 0.9);
\draw[->, thick, DarkPurple!70] (0.7,1.2) -- (1.2, 1.5);
\draw[->, thick, DarkBlue!50] (2.5,0.9) -- (3.1, 0.7);
\end{tikzpicture}\\[1em]

\begin{tcolorbox}[comparisonbox, width=10cm]
\centering\small
\textbf{Prof. José Mauricio Rodas Rodríguez}\\
Universidad de Caldas --- Departamento de Química\\
\textit{Introducción a la Química Computacional (173G7G)}
\end{tcolorbox}
\end{center}
\end{frame}

%==============================================================================
\end{document}
%==============================================================================
